\begin{abstract}
\noindent
Este trabalho apresenta o GraphCube Lab, uma plataforma aberta, interativa e reprodutível para o estudo do cubo mágico com ferramentas de teoria dos grupos, teoria dos grafos e busca heurística. Partindo de uma análise do projeto GraphCube, mostramos que sua representação visual é um diagrama de incidência entre posições de adesivos e camadas --- e não um grafo de estados nem um grafo de Cayley --- e distinguimos formalmente essas três estruturas. O núcleo matemático deriva todas as permutações de uma única descrição geométrica dos adesivos, adota uma convenção de composição explícita e é validado contra resultados conhecidos: o algoritmo de Schreier--Sims reproduz $|G| = \num{43252003274489856000}$ e as ordens de seis subgrupos, e a busca em largura completa do \cube{2} reproduz a distribuição de distâncias publicada nas métricas HTM e QTM. Implementamos BFS, BFS bidirecional, A$^*$ e IDA$^*$ sobre grafos implícitos, heurísticas admissíveis e consistentes obtidas por projeções (\textit{pattern databases}) e o método de duas fases de Kociemba, sempre distinguindo algoritmos que garantem caminhos mínimos dos que produzem apenas soluções válidas. Experimentos com sementes registradas mostram, entre outros resultados, que o comprimento do embaralhamento é um mau substituto para a distância real (que satura em torno da média uniforme \gclMeanDistHTM{}), que a busca bidirecional reduz o crescimento por camada de \gclEoneGrowthBFS{} para \gclEoneGrowthBi{}, que a combinação de duas \textit{pattern databases} reduz os nós expandidos por IDA$^*$ em \gclEthreeMaxOverPerm{}$\times$ em relação à melhor delas isolada, e quantificam a estrutura de grafos de Cayley de \gclEfourSets{} conjuntos de geradores. Todas as visualizações parciais são rotuladas como tais e as animações de busca reproduzem o registro real de execução.

\medskip\noindent\textbf{Palavras-chave:} cubo mágico; teoria dos grupos; grafos de Cayley; algoritmos de busca; heurísticas admissíveis; visualização científica; reprodutibilidade.
\end{abstract}

\begin{otherlanguage}{english}
\begin{abstract}
\noindent
We present GraphCube Lab, an open, interactive and reproducible platform for studying the Rubik's cube through group theory, graph theory and heuristic search. Starting from an analysis of the GraphCube project, we show that its visual representation is a sticker--layer incidence diagram --- neither a state graph nor a Cayley graph --- and formally distinguish these three structures. The mathematical core derives every permutation from a single geometric description of the stickers, uses an explicit composition convention and is validated against known results: the Schreier--Sims algorithm reproduces $|G| = \num{43252003274489856000}$ and the orders of six subgroups, and an exhaustive breadth-first search of the $2\times2\times2$ cube reproduces its published distance distribution in both HTM and QTM. We implement BFS, bidirectional BFS, A$^*$ and IDA$^*$ over implicit graphs, admissible and consistent heuristics obtained from projections (pattern databases) and Kociemba's two-phase method, always separating algorithms that guarantee shortest paths from those that only return valid solutions. Seeded experiments show, among other results, that scramble length is a poor proxy for the true distance, that bidirectional search lowers the per-layer growth rate from \gclEoneGrowthBFS{} to \gclEoneGrowthBi{}, that combining two pattern databases reduces IDA$^*$ expansions by a factor of \gclEthreeMaxOverPerm{} over the best single one, and quantify the structure of Cayley graphs for \gclEfourSets{} generating sets. Partial visualisations are labelled as such and search animations replay the actual execution log.

\medskip\noindent\textbf{Keywords:} Rubik's cube; group theory; Cayley graphs; search algorithms; admissible heuristics; scientific visualisation; reproducibility.
\end{abstract}
\end{otherlanguage}
